% ===========================================================================
% §2 — Le modèle : définition, situations, entrées, hypothèses
% ===========================================================================
\section{Le modèle : périmètre et entrées}
\label{sec:modele}

\subsection{Définition du revenu disponible}

Le \textbf{revenu disponible} (RD) d'un ménage est ce qui lui reste pour
consommer et épargner une fois le système socio-fiscal traversé :

\begin{equation}
\label{eq:rd}
RD \;=\; \underbrace{R}_{\substack{\text{revenu de travail}\\\text{ou de retraite}}}
\;-\; \underbrace{C}_{\text{cotisations}}
\;+\; \underbrace{T_{QC}}_{\substack{\text{transferts}\\\text{québécois}}}
\;-\; \underbrace{I_{QC}}_{\substack{\text{impôt}\\\text{québécois}}}
\;+\; \underbrace{T_{CA}}_{\substack{\text{transferts}\\\text{fédéraux}}}
\;-\; \underbrace{I_{CA}}_{\substack{\text{impôt}\\\text{fédéral}}}
\;-\; \underbrace{G}_{\substack{\text{frais de garde}\\\text{payés}}}
\end{equation}

Tous les montants sont \textbf{annuels} et exprimés en dollars courants de
l'année simulée (2025 ou 2026). Le modèle calcule chaque terme de
l'équation~\eqref{eq:rd} à partir d'un petit nombre d'entrées décrivant le
ménage, puis présente la décomposition complète.

\subsection{Les cinq situations types}

Le modèle reprend les cinq situations du calculateur officiel du MFQ. Chacune
fixe le nombre d'adultes et la nature du revenu saisi :

\begin{table}[ht]
\centering
\small
\begin{tabular}{l c c l}
\toprule
Situation & Adultes & Enfants & Revenu saisi \\
\midrule
Personne vivant seule        & 1 & ---        & travail \\
Famille monoparentale        & 1 & 1 à 5      & travail \\
Couple                       & 2 & 0 à 5      & travail (deux revenus) \\
Retraité vivant seul         & 1 & ---        & retraite (pension privée) \\
Couple de retraités          & 2 & ---        & retraite (deux pensions) \\
\bottomrule
\end{tabular}
\caption{Les cinq situations du modèle. Pour les ménages \og retraités \fg{},
le revenu saisi est un revenu de pension privée (régime d'employeur, FERR,
rente), auquel le modèle ajoute lui-même la pension fédérale de la Sécurité
de la vieillesse et ses suppléments (poste~17).}
\label{tab:situations}
\end{table}

\subsection{Les entrées}

\begin{itemize}
  \item \textbf{Revenu brut} de chaque adulte (travail ou retraite selon la
        situation), de 0 à \D{200000} dans l'interface ;
  \item \textbf{Âge} de chaque adulte --- il commande l'accès à plusieurs
        programmes (aide sociale ajustée à 58 ans, allocation-logement à
        50 ans, Sécurité de la vieillesse à 65 ans, soutien aux aînés à
        70 ans, supplément de PSV à 75 ans) ;
  \item \textbf{Enfants} (0 à 5) : pour chacun, l'âge, les frais de garde
        annuels payés (plafonnés à \D{15000} par enfant dans l'interface) et
        le type de service de garde --- \emph{subventionné} (la place à
        contribution réduite) ou \emph{non subventionné}. Seuls les frais non
        subventionnés ouvrent droit au crédit d'impôt québécois (poste~6) ;
        tous les frais comptent pour la déduction fédérale et pour le coût
        $G$ de l'équation~\eqref{eq:rd}.
\end{itemize}

\subsection{Hypothèses principales}

Le modèle hérite des hypothèses du calculateur officiel ; les connaître évite
de mauvaises interprétations.

\begin{itemize}
  \item \textbf{Une seule source de revenu par adulte} : du travail salarié
        pour les ménages actifs, une pension privée pour les retraités.
        Aucun revenu de placement, d'entreprise, de gain en capital ni de
        REER n'est modélisé ;
  \item \textbf{Aucune déduction facultative} (REER, frais financiers,
        etc.) : seules les déductions automatiques du système sont
        appliquées (déduction pour travailleur au Québec, cotisations RRQ
        supplémentaires, frais de garde au fédéral) ;
  \item \textbf{Enfants de moins de 18 ans}, à charge ;
  \item \textbf{Ménage locataire type} pour l'allocation-logement : le loyer
        n'étant pas saisi, le modèle l'impute selon la composition du ménage
        (poste~10) ;
  \item \textbf{Prestations fiscales versées de juillet à juin} (Allocation
        canadienne pour enfants, crédit TPS, crédit pour la solidarité) :
        le modèle assimile l'année de versement débutant en juillet de
        l'année $Y$ à l'année civile $Y$ ;
  \item \textbf{Couverture RAMQ publique} : les adultes sont présumés
        inscrits au régime public d'assurance médicaments (pas d'assurance
        collective privée).
\end{itemize}
